\begin{topic}[id=eta_aids,
             difficulty={Mittel},
             type={Systemanalyse + Anwendungsbewertung},
             prerequisites={Grundlagen Sensorik; Interesse an Assistenzsystemen},
             practice={optional},
             owner={Dozententeam},
             updated={2026-04-09},
             tags={ ETA, Assistive Tech, Sensorik, Haptik, Barrierefreiheit }]{Electronic Travel Aids}

\TopicDescription{Elektronische Mobilitätshilfen unterstützen Menschen mit Sehbeeinträchtigung. Wir vergleichen Sensorik (Ultraschall, LiDAR, Kamera), Feedbackkanäle (haptisch, auditiv), Algorithmen und Usability-Aspekte. Ziel ist ein praxisnaher Überblick zu Technik, Grenzen und ethischen Fragen.}

\TopicLeitfragen{\item Welche Sensor-/Feedback-Kombinationen sind praxistauglich?
\item Wie adressiert man Fehlalarme und Vertrauen?
\item Wo liegen Datenschutz-/Akzeptanz-Barrieren?
}
\TopicScope{\item Keine Medizinprodukte-Zertifizierung.
\item Keine umfangreiche Bildverarbeitung.
}
\TopicPractice{Optionale Praxisidee: Prototyp-Skizze (Sensor + haptisches Feedback) und Usability-Kriterien.}

\TopicKeywords{ETA, Assistive Tech, Sensorik, Haptik, Barrierefreiheit}

\nocite{etaSurvey}
\end{topic}
